\chapter{Introdução}
\label{cap:introducao}

A ordenação é uma das operações mais estudadas da Ciência da Computação, tanto por
sua ubiquidade prática quanto por seu papel didático na análise de algoritmos.
Segundo \citeonline{knuth1998}, estima-se que uma fração expressiva de todo o tempo
de processamento gasto em sistemas computacionais seja dedicada a ordenar dados ou
a manter dados ordenados. Ordenar não é apenas um fim em si: é a etapa que viabiliza
busca binária, junções eficientes em bancos de dados, algoritmos de geometria
computacional e compressão de dados.

Entre os algoritmos de ordenação interna, o \textit{Quicksort}, proposto por
\citeonline{hoare1962}, ocupa posição de destaque. Seu tempo médio de execução é
$O(n \log n)$ com constantes multiplicativas pequenas, e ele opera essencialmente
\textit{in place}, exigindo memória auxiliar apenas para a pilha de recursão.
Por esse motivo é o algoritmo adotado, em variantes refinadas, pelas bibliotecas
padrão de praticamente todas as linguagens de programação de uso geral.

O Quicksort, contudo, apresenta duas fragilidades bem conhecidas. A primeira é seu
pior caso: quando o pivô escolhido é sistematicamente o menor ou o maior elemento do
subvetor, as partições tornam-se maximamente desbalanceadas e o algoritmo degenera
para $\Theta(n^2)$ comparações. A segunda é de natureza prática: para subvetores
pequenos, o custo fixo da recursão e da partição supera o benefício da estratégia de
divisão e conquista, e algoritmos assintoticamente piores --- como o
\textit{Insertion Sort} --- tornam-se efetivamente mais rápidos.

\section{Problema}

Este trabalho trata do seguinte problema: \textbf{quanto se ganha, na prática, ao
corrigir essas duas fragilidades do Quicksort, e sob quais condições cada correção
compensa?} A pergunta não é respondida apenas pela análise assintótica, uma vez que
ambas as correções preservam a complexidade $O(n \log n)$ do caso médio e atuam
justamente sobre as constantes multiplicativas que a notação assintótica descarta.

\section{Objetivos}

O objetivo geral é realizar um estudo comparativo experimental de três
implementações do Quicksort, conforme o exercício 15 do capítulo 4 de
\citeonline{ziviani2018}:

\begin{enumerate}
  \item \textbf{Quicksort recursivo}: versão clássica, com recursão até subvetores
        unitários;
  \item \textbf{Quicksort híbrido}: versão em que a partição é interrompida para
        subvetores com menos de $M$ elementos, os quais são ordenados por Insertion
        Sort, sendo o valor de $M$ determinado empiricamente;
  \item \textbf{Quicksort híbrido com mediana-de-três}: aprimoramento da versão
        anterior, em que o pivô é a mediana entre o primeiro, o central e o último
        elemento do subvetor.
\end{enumerate}

São objetivos específicos: (i) determinar empiricamente o valor de $M$; (ii) medir
tempo de execução, número de comparações e número de trocas sobre massas de teste
variadas; (iii) forçar explicitamente o pior caso do Quicksort e verificar
experimentalmente seu comportamento quadrático; e (iv) analisar criticamente os
resultados à luz da teoria, justificando as diferenças observadas entre as versões.

\section{Organização do trabalho}

O Capítulo \ref{cap:referencial} apresenta o referencial teórico, com a descrição
dos algoritmos Insertion Sort e Quicksort e a análise de complexidade de ambos,
com destaque para o pior caso do Quicksort. O Capítulo \ref{cap:metodologia}
descreve a metodologia adotada, incluindo as massas de teste, as convenções de
contagem e os cuidados tomados na medição de tempo. O Capítulo
\ref{cap:resultados} apresenta e analisa os resultados obtidos. O Capítulo
\ref{cap:conclusao} sintetiza as conclusões do trabalho.
